\begin{textsample}{NEG Dim 3 – al – Score -1.00 – al/al\_ef\_15.txt}  \label{ex:f3_neg_074}
4.
CONTEXTUALIZANDO AS COMPETÊNCIAS 4.1 Competências de LP dialogando com as competências gerais Embora voltadas ao ensino de Língua Portuguesa, as propostas aqui apresentadas também dizem respeito à articulação do ensino do componente com os demais componentes da área de Linguagens, e também com as outras áreas de conhecimento, reafirmando o caráter integral do ensino no cumprimento de sua função social.
Os conhecimentos da Língua Portuguesa, desenvolvidos a partir de atividades voltadas para o uso e reflexão dos diferentes eixos que organizam o componente, devem estar em serviço da ampliação da participação do estudante em práticas de linguagem de diferentes esferas de atividade social.
As competências dos estudantes, portanto, direitos de aprendizagem e deveres de ensino devem, contemplar o indivíduo de forma contextualizada ao mundo atual, cada vez mais volátil, incerto, complexo e ambíguo.
Para o uso e a reflexão da linguagem desde o primeiro ano do Ensino Fundamental, é necessário garantir, aos estudantes o acesso à diversidade textual e de situações comunicativas da linguagem.
O trabalho com gêneros textuais, respeitando o desenvolvimento cognitivo de cada série/ano, é abordado de forma mais aprofundada em nosso organizador curricular, partindo dos Campos de Atuação até os Desdobramentos Didático-Pedagógicos ( DESDP ).
É importante, também, garantir a oportunidade de reflexão cotidiana sobre as características e o funcionamento da língua.
Isso permitirá que se ampliem os seus repertórios de conhecimento e processo de letramento.
Merece destaque, em função do uso social da língua, o estudo respeitoso e reflexivo das variedades linguísticas, conforme competência específica do componente : Esta 4ª Competência Específica de \textbf{LP}, vai se relacionar diretamente com a 9ª competência geral trazida pela BNCC : O ensino da língua, com foco na concepção enunciativo-discursiva, conforme proposto pelo documento da BNCC, direciona -se também a este documento, visando garantir a reflexão da língua nessa mesma perspectiva.
O que quer dizer que a concepção de reflexão sobre a linguagem visa entender seu funcionamento, nas mais diversas manifestações.
Disso depende o despertar da consciência linguística no estudante ( PEREIRA, 2010 ) e, para tanto, o professor também deve ter esse conhecimento, inclusive nos Anos Iniciais e durante o processo de ortografização.

% matched lemmas: lp
\end{textsample}
